\section{Additional Institutional Details and Fieldwork Information}

\subsection{Ministerial Formula to Calculate High School GPA}\label{sec:ptjeranking}

The Ministry ranks students within each school according to an adjusted high school GPA score, called {\itshape Puntaje Ranking de Notas}, calculated using a specific formula. Here, we report an English translation of official information on the formula, which can be found at \url{https://demre.cl/paes/factores-seleccion/puntaje-ranking}. Although the adjusted formula calculates each student's relative position compared to prior cohorts, the final ranking determining the within-school top 15\% cutoff for PACE seats is computed exclusively among students within the same cohort and school.\par 



First, each student's average grade in each high school year, $GPA_{ig}$, where $g=9,10,11,12$, is rescaled to range from 100 to 1000.\footnote{The Department of Educational Assessment, Measurement, and Registration (DEMRE) publishes tables to rescale GPA, see for example \url{https://demre.cl/proceso-admision/factores-seleccion/tabla-transformacion-nem-5-procesos-grupo-c.php}.} It is then transformed into a relative score, $R_{ig}$, broadly capturing where the student ranks compared to the three prior cohorts of students who completed the same grade in the same school in which the student attended that grade (called the ``reference population"). Letting $\overline{GPA}_g$ denote the average grade in the reference population, and $maxGPA_g$ the highest grade in the reference population, this grade-specific relative ranking is computed according to the following formula:

$$R_{ig}=\overline{GPA}_g+(GPA_{ig}-\overline{GPA}_g)\frac{1000-GPA_{ig}}{maxGPA_g-\overline{GPA}_g}$$

\noindent where $\overline{GPA}_g$ and $maxGPA_g$ are expressed on the same 100 to 1000 scale as $GPA_{ig}$. The final score is then obtained as the following average:

$$adjGPA_{i}=\frac{R_{i1}+R_{i2}+R_{i3}+R_{i4}}{4}.$$

\noindent In our sample, the Pearson correlation coefficient between $adjGPA_{i}$ and the average of raw GPA in the four high school years is $97.77\%$.\par 

The top 15\% cutoff for PACE seats is determined by ranking students in the same school and cohort according to $adjGPA_i$. Within each school and cohort, the top 15\% cutoff is determined as the $85^{th}$ percentile of $adjGPA_i$.

\subsection{PACE Process for the Allocation of Preferential Seats}\label{sec:PACEadmissionprocess}
The PACE and regular applications must be submitted to the centralized system during the same time window. For the cohort included in this study, the PACE and regular admission processes were separate. Students could submit a PACE application list and, separately, a regular application list. In each list a student could include up to ten programs (i.e., college and major combinations), potentially different between the two lists. The two processes were entirely independent, and a student could obtain two admissions simultaneously, one from each process. Here we describe the PACE application and admission process.\par 
For each program listed in their PACE applications, applicants receive a distinct application score, called {\itshape Puntaje de Postulaci\'{o}n PACE} (PPP). The score is calculated based on the applicant's GPA during the four years of high school and attendance during the $11^{th}$ and $12^{th}$ grades. To reduce the occurrence of identical scores across applicants, the score is adjusted for each program, taking into account the program's geographic location and its positional ranking within the applicant's list of preferences.\footnote{The formula is $PPP=(0.8 * PRN + 0.2 * GPA)\cdot (1+bonus_{attendance}+bonus_{geog})+bonus_{listrank}$, where $PRN$ is the {\itshape Puntaje Ranking de Notas} (PRN) used to identify the top $15\%$ of students, which is based on the high school GPA with some adjustments (see Appendix \ref{sec:ptjeranking}), and $GPA$ is the raw high school GPA. The correlation coefficient between the raw GPA and the PRN is around 98\%.  The bonus for attendance rewards high school attendance and it reaches a maximum of 5\% for students who did not miss a single day and drops to 0 for those missing 15\% or more days. A bonus for geographic location of 3.5\% (5\%) is awarded for applications to universities in the same area of Chile (North, Center, or South) (region of Chile) as the student's high school, and the bonus for the rank of the program within the applicant list decreases with the program's rank; it is measured in score points, and it is 25 for applications listed first, typically representing less than 5\% of the total score, and 0 for applications listed tenth. \label{foot:PPPformula} }\par 

Applicants to each program are ranked in descending order based on their application score, and available PACE slots are allocated according to this sequence. Should the number of applicants exceed the available slots, those not immediately admitted are placed on a waiting list for their first-choice program. Subsequently, these candidates are considered for admission to the programs listed as their second choice, following the same order-based allocation process. This procedure is iteratively applied to applicants' subsequent choices. Once an applicant is accepted into a program, they are automatically withdrawn from consideration for any programs ranked lower on their preference list. This measure ensures that no applicant is admitted to more than one program. However, applicants remain eligible for programs ranked higher than their successful application, should they be initially placed on a waiting list for such programs. In instances where a student eligible for a guaranteed PACE slot fails to secure admission in any of their listed preferences, the Ministry of Education employs a proprietary algorithm to determine their placement.\par 

\subsection{Fieldwork Information}\label{sec:fieldwork_info}


 The Ministry of Education encouraged school principals to participate in our study; all the sampled schools agreed to participate. Our fieldworkers visited the schools several times and were able to survey all students who were present. \par
We designed and piloted the surveys. The achievement test questions were developed by the professional testing agencies Aptus Chile and Puntaje Nacional; we extensively piloted the test. \par 
Students filled out paper questionnaires. Schools allowed us to administer our survey during class time. Our survey displaced one lecture. It took students approximately 50 minutes to fill out the questionnaire. At the start of the data collection, fieldworkers explained to students that they would take an achievement test for the first 20 minutes, and that they would be entered into a lottery to win an iPad, with the number of lottery tickets determined by the number of correct answers. At the 20-minute mark, fieldworkers told students to stop working on the achievement test and to proceed to the survey part of the questionnaire. If a student completed the achievement test before the 20 minutes were up, she was allowed to proceed to the survey. \par
To limit the influence of fieldworkers, the instructions were printed on the first page of the survey and the fieldworkers read them aloud. To further harmonize the data collection across fieldworkers, they had to submit checklists to their supervisors. During the first 20 minutes, the fieldworkers acted as invigilators. To further avoid cheating, we produced 6 versions of the achievement test. Versions differed in the question order. To ensure that all students faced questions of increasing difficulty, we assigned questions to three different difficulty categories (based on the difficulty index provided by the testing agencies and on extensive piloting on our target population), and we randomized the order of the questions within each category. Students were told, at the start of the test, that they would not all have identical tests.\par

The questionnaires did not show logos of any Ministry or public agency. 

